\section{疑问与逐节评论}\label{sec:comments}

\subsection{总体问题}
\subsubsection*{Q-001：原文证明是否正确}\label{question:correctness}
\paragraph{日期与来源。}2026-10-13；用户在此前项目交流中明确提出。
\paragraph{问题。}这篇文章的证明是否正确？
\paragraph{原文位置。}主定理 \srclabel{thm:main} 及其全部证明依赖。
\paragraph{已有材料。}第\ref{sec:proof-audit}节保留先前助手的关键步骤审读，第\ref{sec:assessment}节区分助手暂定判断与用户评价。
\recordstatus{待继续逐步讲解和核查；用户尚未给出自己的整体正确性评价。}

\subsubsection*{Q-002：证明方法能否推广至 KLV 多项式}\label{question:klv}
\paragraph{日期与来源。}2026-10-13；用户在此前项目交流中明确提出。
\paragraph{问题。}能否把该证明扩展到 KLV 多项式？
\paragraph{定位。}相关研究问题，不是原文主定理的一部分。
\paragraph{已有材料。}附录\ref{app:klv}保留先前助手的分析、条件命题和有限计算记录。
\recordstatus{待研究；当前未建立一般 KLV 推广，也不能据此判定原文存在缺陷。}

\subsection{引言、定理表述与研究范围}
截至 2026-10-13 当前记录，用户尚未就这一部分表达具体肯定、批评或澄清要求。后续意见按原文命题定位登记，不由助手代拟成用户立场。

\subsection{背景、实现与分次约定}
用户已确认的阅读起点见第\ref{sec:reader-background}节；这是讲解背景，并非对原文的批评。尚未登记针对 \srclabel{bg:normalization} 等具体步骤的用户疑问。

\subsection{二面体组合、反射图与区间搬运}
尚未登记用户对该部分具体引理的疑问或评价。第\ref{audit:dihedral}、\ref{audit:counts}小节的核查内容属于助手分析，不应自动复制为用户提出的问题。

\subsection{根平面局部化与层分解}
尚未登记用户对该部分具体引理的疑问或评价。后续涉及模型、限制映射、局部化或自由性的讨论，应定位到相应引理，并分别记录解释需求与数学质疑。

\subsection{单边比较、互反与同时归纳}
尚未登记用户对该部分具体引理的疑问或评价。第\ref{audit:edge-bound}、\ref{audit:comparison}小节保留的是助手先前的分析，不代表用户已理解或认可。

\subsection{表述、引用与可读性}
用户要求逐步讲解，但尚未对原文的表述质量、引用充分性或可读性作具体评价。不能把逐步讲解的要求自行转换成“论文写得不清楚”。

% 后续具体意见在以上相应小节中追加，以稳定编号和 \label 定位。
% 每条包含：日期、来源、原文位置、用户原话或忠实转述、助手解释、
% 实际核查证据、处理状态及后续更新；不得自行标记用户已理解。
